% Source-derived chapter 03: Equivalences between definitions of the polar multiplicity
% Original source: bounds-stalks-perverse-sheaves-characteristic-p-shende
\chapter{Equivalences between definitions of the polar multiplicity}
\label{bounds-stalks-perverse-sheaves-characteristic-p-shende-chapter-03}
\source{bounds-stalks-perverse-sheaves-characteristic-p-shende}

\begin{remark}[Source section]
\label{bounds-stalks-perverse-sheaves-characteristic-p-shende-chapter-03-source}
\source{bounds-stalks-perverse-sheaves-characteristic-p-shende}
\source{bounds-stalks-perverse-sheaves-characteristic-p-shende}
This chapter reproduces the corresponding section of the retrieved
LaTeX source, including its definitions, statements, examples, and
proofs. It has not yet been translated into Lean.
\notready
\end{remark}

% The source section title was: Equivalences between definitions of the polar multiplicity
\label{equivalences}
\source{bounds-stalks-perverse-sheaves-characteristic-p-shende}

In this section we give an alternate definition of the polar multiplicity, check that it is equivalent to the previous one, and check that both are well-defined.


\begin{defi}
\source{bounds-stalks-perverse-sheaves-characteristic-p-shende}
\notready Let $Y$ be a smooth variety with a map $f$ to a variety $X$ (which may be the identity), and let $x$ be a point on $X$. Let $C_1, C_2$ be algebraic cycles on $Y$ of total dimension $\dim Y$ such that $C_1 \cap C_2 \cap f^{-1}(x)$ is proper. Assume that all connected components of $C_1 \cap C_2$ are either contained in $f^{-1}(x)$ and proper or disjoint from $X$. We define their intersection number locally at $x$  \[ (C_1,C_2)_{Y, x} \] to be the sum of the degrees of the refined intersection $C_1 \cdot C_2$ \cite[p. 131]{Fulton} on all connected components of $C_1 \cap C_2$ contained in $f^{-1}(x)$. \end{defi} 


\begin{lemma}
\source{bounds-stalks-perverse-sheaves-characteristic-p-shende}
\notready\label{intersection-comparison} Let $X$ be a smooth variety. Let $C$ be a conical cycle on the cotangent bundle of $X$ of dimension $\dim X$ and let $x$ be a point on $X$.  Let $\mathbb P(C) \subseteq \mathbb P (T^* X)$ be the projectivization of $C$ inside the projectivization of the cotangent bundle of $X$. Let $i$ be a natural number with $0 \leq i < \dim X$.
\source{bounds-stalks-perverse-sheaves-characteristic-p-shende}

Consider $Y \subset X$ a smooth variety of dimension $\dim X-i$ through $x$ and $V$ a sub-bundle of $T^* X$  of rank $i+1$ on $Y$.  Let $\mathbb P(V) \subseteq \mathbb P( T^* X)$ be the projectivization of $V$ over $Y$. For any $(Y,V)$ such that the strict transforms of $\mathbb P(C)$ and $\mathbb P(V)$ in the blowup of $\mathbb P(T^* X)$ at the fiber over $X$ do not intersect inside the exceptional divisor, the contribution of the fiber over $x$ to the intersection number $\mathbb P(C) \cap \mathbb P(V)$ depends only on $i$ and is independent of $Y,V$.

Furthermore, to satisfy the condition on the strict transform, it is sufficient that the tangent space of $Y$ at $x$ and the fiber of $V$ over $x$ are independent generic subspaces of the tangent and cotangent spaces of $X$ at $x$ respectively. In particular, such a $Y$ and $V$ exist.\end{lemma}

\begin{proof} This is a local question, and we may work locally.  Then given $(Y,V)$ and $(Y',V')$ both satisfying this condition, we may deform one into the other by a connected family of varieties. For instance we may represent $Y$ and $Y'$ as local complete intersections and deform the equations defining $Y$ into the equations defining $Y'$ by convex combination, and similarly for the vector subbundles defining $Y'$ and $V'$. The condition that the intersection of the strict transforms vanishes is an open condition, because the strict transform of $\mathbb P(V)$ varies properly with $Y$ and $V$, so we may assume that there is a family connecting $(Y,V)$ to $(Y',V')$ where every member satisfies this condition. Then because the intersection locus in the blow-up is closed, its image inside $X$ is too, and because it is disjoint from $x$, there must be some neighborhood of $X$ that it doesn't intersect. Then for any $Y_t,V_t$ in the family, the intersection of $\mathbb P(C)$ and $\mathbb P(V_t)$ in $\mathbb P (T^* X)$ is empty in that neighborhood minus $x$, so the contribution to the intersection coming from the fiber over $x$ is constant in the family, and thus is equal for $(Y,V)$ and $(Y',V')$.

For the claim about generic subspaces, note that $C$ has dimension $\dim X$, so $\mathbb P(C)$ has dimension $\dim X-1$, and the intersection of its strict transform with the fiber has dimension $\dim X-1$.  The fiber of the blowup is isomorphic to $\mathbb P( (TX)_x ) \times \mathbb P( (T^*X)_x )$, of dimension $2 \dim X - 2$, and the strict transform of $\mathbb P(V)$ is $\mathbb P((TY)_x) \times \mathbb P ( V_x)$, of dimension $\dim X - i -1 + i = \dim X-1$. If we take $(TY)_x$ and $V_x$ to be general subspaces, this intersection will have the expected dimension, which is $-1$, and hence be empty. \end{proof}



\begin{defi}
\source{bounds-stalks-perverse-sheaves-characteristic-p-shende}
\notready\label{polar-multiplicity-2} Let $X$ be a smooth variety. Let $C$ be a conical cycle on the cotangent bundle $T^* X$ of $X$ of dimension $\dim X$ and let $x$ be a point on $X$.
\source{bounds-stalks-perverse-sheaves-characteristic-p-shende}

For $0 \leq i< \dim X$, let $Y$ be a sufficiently general smooth subvariety of $X$ of codimension $i$ passing through $x$ and let $V$ be a sufficiently general sub-bundle of $T^* X$ over $Y$ with rank $i+1$. Define the $i$th polar multiplicity of $C$ at $x$ to be the intersection number \[ ( \mathbb P(C) , \mathbb P(V))_{ \mathbb P(T^* X), x} \] where $\mathbb P(T^* X)$ is the projectivization of the vector bundle $T^* X$.

Here ``sufficiently general" means that the strict transform of $\mathbb P(V)$ in the blowup of $\mathbb P(T^* X)$ at the fiber over $x$ does not intersect the strict transform of $\mathbb P(C)$ in that same blowup within the fiber over $x$.

For $i=\dim X$, define the $i$th polar multiplicity of $C$ at $x$ to be the multiplicity of the zero section in $C$. \end{defi}

It follows from Lemma \ref{intersection-comparison} that this is well-defined.




\begin{lemma}
\source{bounds-stalks-perverse-sheaves-characteristic-p-shende}
\notready\label{multiplicity-polar} Definitions \ref{polar-multiplicity} and \ref{polar-multiplicity-2} are equivalent.
\source{bounds-stalks-perverse-sheaves-characteristic-p-shende}
 \end{lemma}

\begin{proof} By definition, the multiplicity of an algebraic cycle at a point is the local intersection number with a sufficiently general smooth scheme passing through that point. For $Y$ a sufficiently general smooth subscheme of $X$ of dimension $n-i$, we have an identity of intersection numbers \[ (\pi_* ( \mathbb P(C) \cap \mathbb P(V) , Y)_{X,x}) = (\mathbb P(C) \cap \mathbb P(V), \pi^* Y)_{\mathbb P(T^*X), X} =  (\mathbb P(C), \mathbb P(V) \cap \pi^* Y)_{\mathbb P(T^*X), X} .\]

Note that $\mathbb P(V) \cap \pi^* Y$ is simply the projectivization of the restriction $V'$ of $V$ to $Y$. So to check that this is the polar multiplicity, it suffices to check that if $V_x$ is sufficiently general, and $Y$ is sufficiently general depending on $V$, that the restriction of $V$ to $Y$ is sufficiently general in the sense of Definition \ref{polar-multiplicity-2}. This occurs when the intersection of the strict transform of $\mathbb P(C)$ with the strict transform of $\mathbb P(V')$ in the exceptional divisor of the blowup of $\mathbb P (T^* X)$ at the fiber over $x$ vanishes.

The exceptional divisor is isomorphic to $\mathbb P ( (TX)_x ) \times \mathbb P( (T^*X)_x) $. Inside it, the strict transform of $\mathbb P(V')$ is $\mathbb P ( (TY)_x) \times \mathbb P (  V_x)$. The intersection of the strict transform of $\mathbb P(C)$ with the exceptional divisor has dimension at most $\dim \mathbb P(C) -1 =\dim C-2= n-2$. For $V$ of dimension $i+1$, $\mathbb P( (TX)_x) \times \mathbb P( V_x)$ has codimension $n-i-1$, so for $V_x$ sufficiently general,  the intersection of the strict transform with $( \mathbb P( (TX)_x) \times \mathbb P( V_x))$ has dimension $i-1$. Then for general $\mathbb P( (TY)_x)$ of codimension $i$, the intersection of $\mathbb P ( (TY)_x) \times \mathbb P (  V_x)$ with the strict transform is empty.  
\end{proof}
